% ── ادامه بخش مشارکت سیاسی ──────────────────────────────
\infobox{boxgreen}{روند مثبت: کاهش رأی فرقه‌ای}{%
سهم «تعلق فرقه‌ای/قومی» به‌عنوان معیار انتخاب از ۳۰٪ (۲۰۰۸)
به ۲۰٪ (۲۰۱۳) کاهش یافت. در مقابل «صداقت و پاکدستی»
از ۲۲٪ به ۳۰٪ رسید.
اما باور به تأثیر رأی از ۶۵٪ به ۳۸٪ سقوط کرد ─
نشانه‌ای از بی‌اعتمادی ساختاری به کل فرایند انتخاباتی.
این پارادوکس ─ رأی‌دهندگان آگاه‌تر اما ناامیدتر ─
در انتخابات ۲۰۱۸ و ۲۰۲۱ با پایین‌ترین مشارکت تاریخ عراق
(≈۴۰٪) تجلی یافت.
}


% ████████████████████████████████████████████████████████████
%  بخش ۱۲ ─ نظرسنجی‌های محرمانه: D3 Systems
% ████████████████████████████████████████████████████████████
\section{نظرسنجی‌های محرمانه: صدایی که واشنگتن شنید اما نادیده گرفت}

\noindent
بخشی از داده‌های ارزشمند درباره‌ی افکار عمومی عراق، در
نظرسنجی‌هایی محرمانه جمع‌آوری شد که برای مصرف داخلی
دولت آمریکا طراحی شده بودند. انتشار بعدی آن‌ها از طریق
\lr{WikiLeaks} و درخواست‌های \lr{FOIA} تصویری خاص ارائه می‌دهد.

\begin{surveycard}{\lr{D3 Systems / KA Research} ─ برای وزارت خارجه و ارتش آمریکا}
\begin{tabularx}{\textwidth}{>{\bfseries\small}r >{\small}X}
مجری          & \lr{D3 Systems} (وین، ویرجینیا) + \lr{KA Research} (استانبول) \\
سفارش‌دهنده   & وزارت خارجه (\lr{State Dept.}) و وزارت دفاع (\lr{DoD}) آمریکا \\
فراوانی       & تقریباً هر ۳ ماه ─ بیش از ۲۰ موج (۲۰۰۴–۲۰۱۱) \\
نمونه         & \ptho{۲}{۰۰۰}–\ptho{۵}{۰۰۰} \\
ویژگی         & پرسش‌های عملیاتی–نظامی · برخی محرمانه \\
\end{tabularx}
\end{surveycard}

\begin{table}[H]
\centering
\caption{یافته‌های کلیدی نظرسنجی‌های محرمانه‌ی \lr{D3/KA}}
\label{tab:d3-findings}
\small
\begin{tabularx}{\textwidth}{>{\centering\arraybackslash\small}p{1.2cm} >{\bfseries\small}r >{\small}X}
\toprule
\textbf{سال} & \textbf{پرسش} & \textbf{یافته‌ی کلیدی} \\
\midrule
۲۰۰۶ & حملات علیه ائتلاف & ۶۱٪ عراقی‌ها حملات علیه نیروهای آمریکایی را تأیید کردند \\
\addlinespace
۲۰۰۷ & اعتماد به آمریکا & فقط ۱۸٪ به آمریکا اعتماد داشتند (پایین‌ترین سطح) \\
\addlinespace
۲۰۰۸ & آینده‌ی بهتر & ۵۵٪ معتقد بودند زندگی فرزندانشان بهتر خواهد بود \\
\addlinespace
۲۰۰۸ & خروج آمریکا & ۷۳٪ خواهان خروج بر اساس جدول زمانی بودند \\
\addlinespace
۲۰۰۹ & اعتماد به ارتش عراق & ۷۰٪ به ارتش عراق اعتماد داشتند \\
\addlinespace
۲۰۱۰ & اعتماد به مالکی & ۴۸٪ عملکرد مالکی را مثبت ارزیابی کردند \\
\addlinespace
۲۰۱۱ & آمادگی ارتش عراق & ۵۵٪ معتقد بودند ارتش عراق آماده‌ی جایگزینی آمریکاست \\
\bottomrule
\end{tabularx}
\end{table}

\didyouknow{%
  نظرسنجی‌های محرمانه‌ی \lr{D3 Systems} نشان می‌دادند
  حمایت عراقی‌ها از حملات علیه نیروهای آمریکایی از
  ۲۰۰۴ به ۲۰۰۶ \textbf{دو برابر} شده بود ─ اما این
  اطلاعات ماه‌ها پیش از انتشار عمومی در اختیار
  فرماندهان نظامی قرار داشت. سؤال آن است:
  چرا سیاست تغییر نکرد؟
}


% ████████████████████████████████████████████████████████████
%  بخش ۱۳ ─ جدول فراگیر خلاصه
% ████████████████████████████████████████████████████████████
\section{جدول فراگیر: خلاصه‌ی بیست ساله}

\begin{table}[H]
\centering
\caption{جدول فراگیر: پرسش‌های کلیدی و روند پاسخ‌ها (۲۰۰۳–۲۰۲۲)}
\label{tab:mega-summary}
\small
\begin{tabularx}{\textwidth}{>{\bfseries\scriptsize}p{2.6cm} *{8}{>{\centering\arraybackslash\scriptsize}X}}
\toprule
\textbf{پرسش} & \textbf{۲۰۰۳} & \textbf{۲۰۰۴} & \textbf{۲۰۰۵} & \textbf{۲۰۰۷} & \textbf{۲۰۰۹} & \textbf{۲۰۱۳} & \textbf{۲۰۱۹} & \textbf{۲۰۲۲} \\
\midrule
حمله درست بود         & --  & ۴۸٪ & ۴۱٪ & ۳۴٪ & ۶۱٪ & --  & --  & --  \\
\addlinespace
جهت کشور درست         & --  & ۵۱٪ & ۴۴٪ & ۲۲٪ & ۴۱٪ & ۳۰٪ & ۱۵٪ & ۲۲٪ \\
\addlinespace
زندگی بهتر از قبل     & ۵۸٪ & ۵۶٪ & ۴۸٪ & ۲۶٪ & ۴۵٪ & ۳۲٪ & ۲۰٪ & ۲۸٪ \\
\addlinespace
وضع اقتصادی خوب       & --  & ۴۲٪ & ۳۵٪ & ۲۰٪ & ۳۵٪ & ۲۵٪ & ۱۵٪ & ۲۰٪ \\
\addlinespace
امنیت محل خوب         & ۵۵٪ & ۶۳٪ & ۵۵٪ & ۲۶٪ & ۵۸٪ & ۴۵٪ & ۵۵٪ & ۶۰٪ \\
\addlinespace
اعتماد به دولت        & ۵۰٪ & ۵۰٪ & ۴۵٪ & ۲۲٪ & ۳۵٪ & ۱۸٪ & ۱۰٪ & ۱۵٪ \\
\addlinespace
خواهان خروج ائتلاف    & --  & ۵۱٪ & ۶۵٪ & ۷۸٪ & ۶۵٪ & --  & ۷۵٪ & --  \\
\addlinespace
دموکراسی بهترین نظام  & ۷۲٪ & --  & ۵۵٪ & ۵۰٪ & --  & ۵۳٪ & ۵۵٪ & ۵۰٪ \\
\addlinespace
فساد بزرگ‌ترین مشکل   & ۸٪  & ۵٪  & ۵٪  & ۵٪  & ۲۰٪ & ۲۸٪ & ۴۵٪ & ۳۸٪ \\
\bottomrule
\end{tabularx}
\end{table}


% ████████████████████████████████████████████████████████████
%  بخش ۱۴ ─ پنج درس کلان
% ████████████████████████████████████████████████████████████
\section{پنج درس کلان از بیست سال نظرسنجی}

\subsection{درس اول: هویت فرقه‌ای قوی‌ترین پیش‌بین}

\begin{figurebox}
\begin{figure}[H]
\centering
\begin{tikzpicture}
\begin{axis}[
  ybar, width=15cm, height=8cm,
  bar width=11pt,
  xlabel={\rl{پرسش}},
  ylabel={\rl{درصد پاسخ مثبت}},
  symbolic x coords={%
    {\rl{ارزش سرنگونی}},
    {\rl{جهت درست}},
    {\rl{زندگی بهتر}},
    {\rl{امنیت خوب}},
    {\rl{اعتماد دولت}},
    {\rl{حملات مشروع}}},
  xtick=data,
  xticklabel style={rotate=30, anchor=east, font=\small},
  ymin=0, ymax=100,
  legend pos=north east,
  legend style={font=\small},
  grid=major, grid style={dashed, gray!25},
  title={\rl{\textbf{شکاف فرقه‌ای در پاسخ‌ها (۲۰۰۹)}}},
  title style={font=\small},
  nodes near coords,
  every node near coord/.append style={font=\tiny},
]
\addplot[fill=green!60] coordinates {
  ({\rl{ارزش سرنگونی}},85)
  ({\rl{جهت درست}},68)
  ({\rl{زندگی بهتر}},70)
  ({\rl{امنیت خوب}},85)
  ({\rl{اعتماد دولت}},55)
  ({\rl{حملات مشروع}},3)
};
\addlegendentry{\rl{کردها}}
\addplot[fill=blue!60] coordinates {
  ({\rl{ارزش سرنگونی}},75)
  ({\rl{جهت درست}},48)
  ({\rl{زندگی بهتر}},58)
  ({\rl{امنیت خوب}},55)
  ({\rl{اعتماد دولت}},55)
  ({\rl{حملات مشروع}},22)
};
\addlegendentry{\rl{شیعیان}}
\addplot[fill=red!60] coordinates {
  ({\rl{ارزش سرنگونی}},23)
  ({\rl{جهت درست}},18)
  ({\rl{زندگی بهتر}},22)
  ({\rl{امنیت خوب}},40)
  ({\rl{اعتماد دولت}},15)
  ({\rl{حملات مشروع}},65)
};
\addlegendentry{\rl{سنی‌ها}}
\end{axis}
\end{tikzpicture}
\caption{در هر پرسش، هویت فرقه‌ای بزرگ‌ترین متغیر تعیین‌کننده است}
\label{fig:sectarian-gap-2009}
\end{figure}
\end{figurebox}


\subsection{درس دوم: فرسایش خوش‌بینی اولیه}

\begin{table}[H]
\centering
\caption{فرسایش خوش‌بینی: مقایسه‌ی اوج با حضیض}
\label{tab:optimism-erosion}
\begin{tabularx}{\textwidth}{>{\bfseries\small}r >{\small\centering\arraybackslash}p{1.5cm} >{\small\centering\arraybackslash}p{1.5cm} >{\small\centering\arraybackslash}p{1.5cm} >{\small\centering\arraybackslash}p{1.5cm} >{\small\centering\arraybackslash}p{1.5cm}}
\toprule
\textbf{شاخص} & \textbf{اوج} & \textbf{سال} & \textbf{حضیض} & \textbf{سال} & \textbf{افت} \\
\midrule
خوش‌بینی به آینده   & ۶۷٪ & ۲۰۰۳ & ۱۵٪ & ۲۰۱۹ & $-$۵۲ \\
اعتماد به دولت     & ۵۰٪ & ۲۰۰۴ & ۱۰٪ & ۲۰۱۹ & $-$۴۰ \\
جهت کشور درست      & ۵۱٪ & ۲۰۰۴ & ۱۵٪ & ۲۰۱۹ & $-$۳۶ \\
اعتماد به احزاب    & ۳۵٪ & ۲۰۰۴ & ۵٪  & ۲۰۱۹ & $-$۳۰ \\
رضایت از حمله      & ۶۱٪ & ۲۰۰۹ & ۲۲٪ & ۲۰۰۷ & $-$۳۹ \\
\bottomrule
\end{tabularx}
\end{table}


\subsection{درس سوم: پایداری شگفت‌انگیز تقاضا برای دموکراسی}

\infobox{boxgreen}{پارادوکس دموکراسی عراقی}{%
علی‌رغم:
\begin{itemize}[rightmargin=0.5cm]
\item سرخوردگی عمیق از عملکرد نظام سیاسی
\item فروپاشی اعتماد به نهادهای حکومتی
\item تجربه‌ی فاجعه‌بار جنگ فرقه‌ای و داعش
\end{itemize}
حمایت عراقی‌ها از \textbf{اصل دموکراسی} هرگز به زیر ۴۸٪ نرفت
(بارومتر عرب).
عراقی‌ها بین \textbf{ایده‌ی دموکراسی} و \textbf{عملکرد آن}
تمایز قائل‌اند.
این سرمایه‌ی اجتماعی نشان می‌دهد زمینه‌ی دموکراسی‌سازی
هنوز وجود دارد ─ اما اعتمادسازی نهادی شرط اول است.
}


\subsection{درس چهارم: شکاف نسلی فزاینده}

\begin{table}[H]
\centering
\caption{شکاف نسلی: جوانان در برابر مسن‌ترها (۲۰۱۹)}
\label{tab:generation-gap}
\begin{tabularx}{\textwidth}{>{\bfseries\small}r >{\centering\arraybackslash\small}X >{\centering\arraybackslash\small}X >{\centering\arraybackslash\small}X}
\toprule
\textbf{شاخص} & \textbf{۱۸–۲۹} & \textbf{۳۰–۴۹} & \textbf{۵۰+} \\
\midrule
اعتماد به دولت          & ۷٪  & ۱۰٪ & ۱۵٪ \\
قصد مهاجرت              & ۴۵٪ & ۲۸٪ & ۱۰٪ \\
فساد بزرگ‌ترین مشکل     & ۵۰٪ & ۴۵٪ & ۳۵٪ \\
دموکراسی بهترین نظام    & ۵۲٪ & ۵۸٪ & ۵۵٪ \\
شرکت در اعتراضات        & ۳۰٪ & ۱۵٪ & ۵٪  \\
نارضایتی از محاصصه      & ۷۵٪ & ۶۵٪ & ۵۰٪ \\
هویت «عراقی» اول         & ۵۵٪ & ۴۵٪ & ۳۵٪ \\
\bottomrule
\end{tabularx}
\end{table}

\infobox{boxblue}{تحلیل نسلی: موتور تغییر آینده}{%
نسل پس از ۲۰۰۳ (متولدان ۱۹۹۰ به بعد):
\begin{itemize}[rightmargin=0.5cm]
\item \textbf{هویت ملی قوی‌تر:} ۵۵٪ خود را ابتدا «عراقی» می‌دانند (نه شیعه/سنی/کرد)
\item \textbf{سرخوردگی عمیق‌تر:} کمترین اعتماد به نهادها در تاریخ نظرسنجی‌ها
\item \textbf{تمایل بالای مهاجرت:} ۴۵٪ ─ فرار مغزها به‌مثابه تهدید وجودی
\item \textbf{رادیکالیسم مدنی:} بیشترین مشارکت در جنبش تشرین
\item \textbf{رد محاصصه:} ۷۵٪ مخالف نظام سهمیه‌بندی فرقه‌ای
\item این نسل ممکن است هم \textbf{بزرگ‌ترین تهدید} (مهاجرت/بی‌ثباتی)
      و هم \textbf{بزرگ‌ترین فرصت} (اصلاح فرافرقه‌ای) عراق باشد
\end{itemize}
}


\subsection{درس پنجم: زنان عراقی ─ صداهای ناشنیده}

\begin{table}[H]
\centering
\caption{تفاوت‌های جنسیتی در نظرسنجی‌ها (میانگین ۲۰۰۹–۲۰۱۹)}
\label{tab:gender-gap}
\begin{tabularx}{\textwidth}{>{\bfseries\small}r >{\centering\arraybackslash\small}X >{\centering\arraybackslash\small}X >{\centering\arraybackslash\small}X}
\toprule
\textbf{شاخص} & \textbf{مرد} & \textbf{زن} & \textbf{فاصله} \\
\midrule
احساس امنیت شبانه          & ۵۲٪ & ۳۵٪ & ۱۷ واحد \\
رضایت از وضعیت اقتصادی    & ۳۵٪ & ۲۸٪ & ۷ واحد \\
مشارکت در انتخابات         & ۶۵٪ & ۵۵٪ & ۱۰ واحد \\
باور به تأثیر رأی          & ۴۵٪ & ۳۵٪ & ۱۰ واحد \\
موافقت با حقوق برابر زنان  & ۳۵٪ & ۵۵٪ & ۲۰ واحد \\
نگرانی از خشونت خانگی     & ۲۵٪ & ۵۵٪ & ۳۰ واحد \\
تمایل به مهاجرت           & ۳۵٪ & ۲۵٪ & ۱۰ واحد \\
\bottomrule
\end{tabularx}
\end{table}

\didyouknow{%
  بزرگ‌ترین شکاف جنسیتی در نظرسنجی‌های عراق مربوط به
  \textbf{نگرانی از خشونت خانگی} است:
  ۵۵٪ زنان در برابر ۲۵٪ مردان ─ فاصله‌ی ۳۰ واحدی.
  این عدد نشان می‌دهد خشونت خانگی نه‌تنها «مسئله‌ی زنان»
  بلکه یکی از عمیق‌ترین شکاف‌های ادراکی جامعه‌ی عراق است.
}


% ████████████████████████████████████████████████████████████
%  بخش ۱۵ ─ فهرست جامع پرسش‌ها
% ████████████████████████████████████████████████████████████
\section{ضمیمه: فهرست جامع پرسش‌های کلیدی}

\begin{longtable}{%
  >{\scriptsize\bfseries}p{1.2cm}
  >{\scriptsize}p{2.2cm}
  >{\scriptsize}p{5.2cm}
  >{\scriptsize}p{2.2cm}
  >{\scriptsize}p{1.8cm}}
\caption{فهرست جامع پرسش‌های کلیدی نظرسنجی‌ها در عراق}
\label{tab:all-questions} \\
\toprule
\textbf{کد} & \textbf{مؤسسه} & \textbf{متن پرسش (انگلیسی)} & \textbf{موضوع} & \textbf{سال‌ها} \\
\midrule
\endfirsthead

\multicolumn{5}{r}{\small\itshape ادامه از صفحه‌ی قبل} \\
\toprule
\textbf{کد} & \textbf{مؤسسه} & \textbf{متن پرسش} & \textbf{موضوع} & \textbf{سال‌ها} \\
\midrule
\endhead

\midrule
\multicolumn{5}{l}{\small\itshape ادامه در صفحه‌ی بعد} \\
\endfoot

\bottomrule
\endlastfoot

% ── الف. ارزیابی حمله ──
\multicolumn{5}{r}{\textbf{الف. ارزیابی حمله و سرنگونی}} \\
\midrule
Q-A1 & \lr{ABC/BBC}   & \lr{Was ousting Saddam worth it?}         & ارزش سرنگونی     & ۰۴–۰۹ \\
Q-A2 & \lr{ABC/BBC}   & \lr{Right direction or wrong direction?}  & جهت کشور         & ۰۴–۰۹ \\
Q-A3 & \lr{ABC/BBC}   & \lr{Life compared to before the war?}    & مقایسه‌ی زندگی   & ۰۴–۰۹ \\
Q-A4 & \lr{Pew}       & \lr{Favorable opinion of the US?}        & نگرش به آمریکا   & ۰۴–۱۵ \\
Q-A5 & \lr{Zogby}     & \lr{Has war brought more democracy?}     & دموکراسی منطقه   & ۰۴–۰۸ \\
\addlinespace

% ── ب. امنیت ──
\multicolumn{5}{r}{\textbf{ب. امنیت}} \\
\midrule
Q-B1 & \lr{ABC/BBC}   & \lr{Security in your area?}              & امنیت محلی       & ۰۴–۰۹ \\
Q-B2 & \lr{Gallup}    & \lr{Feel safe walking at night?}         & امنیت شبانه      & ۰۶–۱۱ \\
Q-B3 & \lr{ABC/BBC}   & \lr{Attacks on coalition acceptable?}    & مشروعیت حملات    & ۰۴–۰۹ \\
Q-B4 & \lr{Gallup}    & \lr{Victim of assault/mugging?}          & تجربه‌ی خشونت    & ۰۷–۱۱ \\
Q-B5 & \lr{ABC/BBC}   & \lr{US presence helps/hurts security?}   & تأثیر حضور آمریکا & ۰۴–۰۹ \\
\addlinespace

% ── پ. دموکراسی و حکمرانی ──
\multicolumn{5}{r}{\textbf{پ. دموکراسی و حکمرانی}} \\
\midrule
Q-C1 & \lr{Arab Bar.} & \lr{Democracy better than other systems?} & مطلوبیت دموکراسی & ۰۶–۲۲ \\
Q-C2 & \lr{Pew}       & \lr{Can democracy work here?}            & کارایی دموکراسی  & ۰۴–۱۲ \\
Q-C3 & \lr{WVS}       & \lr{Strong leader without parliament?}   & رهبر قوی         & ۰۴, ۱۳ \\
Q-C4 & \lr{Arab Bar.} & \lr{Trust in government?}                & اعتماد به دولت   & ۰۶–۲۲ \\
Q-C5 & \lr{Arab Bar.} & \lr{Trust in parliament?}                & اعتماد به پارلمان & ۰۶–۲۲ \\
Q-C6 & \lr{ABC/BBC}   & \lr{Confidence in political parties?}    & اعتماد به احزاب  & ۰۴–۰۹ \\
Q-C7 & \lr{IRI}       & \lr{Rate PM performance?}                & عملکرد نخست‌وزیر  & ۰۴–۱۳ \\
Q-C8 & \lr{IRI}       & \lr{Vote makes a difference?}            & تأثیر رأی        & ۰۵–۱۳ \\
Q-C9 & \lr{NDI}       & \lr{Best government structure?}          & ساختار حکومت     & ۰۴–۰۹ \\
\addlinespace

% ── ت. اقتصاد و خدمات ──
\multicolumn{5}{r}{\textbf{ت. اقتصاد و خدمات}} \\
\midrule
Q-D1 & \lr{Gallup}    & \lr{Economic conditions getting better?} & اقتصاد محلی      & ۰۶–۱۱ \\
Q-D2 & \lr{Gallup}    & \lr{Satisfied with standard of living?}  & رضایت از زندگی   & ۰۶–۱۱ \\
Q-D3 & \lr{IRI}       & \lr{Rate electricity availability?}      & کیفیت برق        & ۰۴–۱۳ \\
Q-D4 & \lr{IRI}       & \lr{Rate water availability?}            & کیفیت آب         & ۰۴–۱۳ \\
Q-D5 & \lr{IRI}       & \lr{Rate healthcare?}                    & بهداشت           & ۰۴–۱۳ \\
Q-D6 & \lr{IRI}       & \lr{Rate employment opportunities?}      & اشتغال           & ۰۴–۱۳ \\
Q-D7 & \lr{Arab Bar.} & \lr{Government performance on corruption?} & مبارزه با فساد & ۱۳–۲۲ \\
\addlinespace

% ── ث. هویت و ارزش‌ها ──
\multicolumn{5}{r}{\textbf{ث. هویت و ارزش‌ها}} \\
\midrule
Q-E1 & \lr{ABC/BBC}   & \lr{Primary identity?}                   & هویت اول         & ۰۷–۰۹ \\
Q-E2 & \lr{NDI}       & \lr{Can groups live together?}           & همزیستی          & ۰۴–۰۹ \\
Q-E3 & \lr{WVS}       & \lr{Importance of religion?}             & اهمیت مذهب       & ۰۴, ۱۳ \\
Q-E4 & \lr{WVS}       & \lr{Men better political leaders?}       & نگرش جنسیتی      & ۰۴, ۱۳ \\
Q-E5 & \lr{Pew}       & \lr{Concern about Sunni-Shia tensions?}  & تنش فرقه‌ای      & ۰۶–۱۲ \\
Q-E6 & \lr{Gallup}    & \lr{Freedom to choose life?}             & آزادی فردی       & ۰۶–۱۱ \\
\addlinespace

% ── ج. اولویت‌ها ──
\multicolumn{5}{r}{\textbf{ج. اولویت‌ها و چالش‌ها}} \\
\midrule
Q-F1 & \lr{ABC/BBC}   & \lr{Biggest problem facing country?}     & اولویت اول       & ۰۴–۰۹ \\
Q-F2 & \lr{Arab Bar.} & \lr{Most important challenge?}           & مهم‌ترین چالش     & ۰۶–۲۲ \\
Q-F3 & \lr{Gallup}    & \lr{Cantril ladder: present?}            & رضایت فعلی       & ۰۶–۱۱ \\
Q-F4 & \lr{Gallup}    & \lr{Cantril ladder: 5 years?}            & امید به آینده    & ۰۶–۱۱ \\

\end{longtable}


% ████████████████████████████████████████████████████████████
%  جمع‌بندی فصل
% ████████████████████████████████████████████████████████████
\begin{tcolorbox}[
  enhanced,
  colback=gray!5,
  colframe=black!70,
  fonttitle=\bfseries\large,
  title={جمع‌بندی فصل: صدای ۲۰۰٬۰۰۰ عراقی},
  arc=3mm,
  boxrule=1pt,
]
\begin{enumerate}[rightmargin=0.5cm, itemsep=4pt]

\item \textbf{خوش‌بینی اولیه فروریخت:}
  از ۶۷٪ امید (۲۰۰۳) به ۱۵٪ (۲۰۱۹) ─ اما هرگز صفر نشد.

\item \textbf{هویت فرقه‌ای تعیین‌کننده‌ی همه‌چیز است:}
  تفاوت ۶۰+ واحدی بین کردها و سنی‌ها در اکثر پرسش‌ها.

\item \textbf{اعتماد به سیاسیون سقوط کرد، اعتماد به ارتش صعود:}
  احزاب ۸٪ · ارتش ۷۲٪ ─ خطر نظامی‌شدن سیاست.

\item \textbf{فساد جایگزین امنیت شد:}
  از ۸٪ (۲۰۰۳) به ۴۵٪ (۲۰۱۹) ─ مردم فساد را ریشه‌ی همه‌ی بحران‌ها می‌دانند.

\item \textbf{دموکراسی هنوز زنده است:}
  حمایت از اصل دموکراسی هرگز زیر ۴۸٪ نرفت ─ سرمایه‌ای برای آینده.

\item \textbf{نسل جوان موتور تغییر است:}
  ۵۵٪ هویت ملی · ۷۵٪ ضد محاصصه · ۴۵٪ خواهان مهاجرت ─
  هم تهدید هم فرصت.

\item \textbf{زنان ناامن‌تر، محروم‌تر و ناشنیده‌ترند:}
  شکاف ۳۰ واحدی در نگرانی از خشونت خانگی.

\end{enumerate}

\vspace{6pt}
\noindent
\textit{%
  این داده‌ها حاصل کار میدانی هزاران پرسشگر عراقی
  در شرایط جنگی است.
  پیام تجمیعی آن‌ها روشن است: عراقی‌ها
  \textbf{دموکراسی می‌خواهند}، \textbf{از فساد خسته‌اند}،
  و \textbf{به سیاسیون‌شان اعتماد ندارند}.
  پرسش آن است که آیا نسل تشرین می‌تواند
  این سرمایه‌ی نارضایتی را به تغییر سازنده تبدیل کند.%
}
\end{tcolorbox}